\hypertarget{inspections-surrendered-rights-and-changed-employee-benefits}{%
\subsection{Inspections, surrendered rights and changed employee benefits}\label{inspections-surrendered-rights-and-changed-employee-benefits}}

A replacement inspection, a smaller leased area and a pension curtailment all remove something from an existing arrangement, but they remove different economic rights or obligations. The accounting must identify what survives before measuring the new arrangement. The following methods use the retained historical EU text and synthetic pretax amounts. Actual authority, contract and actuarial facts remain necessary for application.

\hypertarget{replacing-the-inspection-component}{%
\subsubsection{Replacing the inspection component}\label{replacing-the-inspection-component}}

IAS~16.14 \citep[IAS~16.12--14,43--55]{eu20231803} distinguishes a major inspection required for continued operation from the physical parts inspected. Capitalize its cost as a replacement only if the recognition criteria hold; derecognize the remaining previous inspection cost. That obligation persists even if the original purchase did not separately identify an inspection component. The estimated cost of a comparable future inspection may indicate the original component cost; it is not automatically today's replacement expenditure or an additional historical purchase.

The component record carries asset identity, inspection cohort, original cost allocation and evidence, availability date, residual value, consumption pattern, useful life until the next inspection, accumulated depreciation/impairment and the replacement event. Allocate the initial asset's existing total cost among its components; do not increase that total while recognizing the old inspection. Advance depreciation to the replacement date, remove the old carrying amount, add eligible new cost and begin the new component's depreciation when available. IAS 16.43--55 requires separate significant-component depreciation and retains depreciation during idle periods unless the applicable consumption method gives zero. Ordinary servicing is expense under IAS 16.12, not a new inspection asset. Other production-cost rules may place depreciation in inventory; these examples assume a non-production use, so charges enter profit.

Suppose original equipment cost 120 includes an inspection component 20. The remaining physical part costs 100, has ten-year straight-line life and zero residual; the inspection has four-year life and zero residual. After two full years the physical part carries 80 and the old inspection 10. A qualifying inspection performed early costs 24 and is available immediately for four years. Derecognize the old 10 as loss, capitalize 24, and retain physical carrying amount 80. The equipment now carries 104. One subsequent full year charges 10 physical and 6 inspection depreciation, leaving 88. Expensing all 24, retaining both inspection cohorts, or depreciating 104 uniformly over the physical life would compute different targets.

With opening cash 100 and equipment 90, equity 190, replacement leaves cash 76, equipment 104 and equity 180. Subsequent depreciation 16 leaves equity 164 and equipment 88. An otherwise separate routine service costing 3 reduces cash and profit by 3 and leaves equipment unchanged. An unavailable old-cost allocation does not justify writing off zero: reconstruct a supported component estimate or retain the unresolved derecognition amount. Revaluation/impairment histories must first update the component carrying amounts under their own methods.

\hypertarget{reducing-a-lease-and-paying-usage-dependent-rent}{%
\subsubsection{Reducing a lease and paying usage-dependent rent}\label{reducing-a-lease-and-paying-usage-dependent-rent}}

IFRS~16.44 \citep[IFRS~16.38--46]{eu20231803} first tests whether an added right and commensurate consideration make a separate lease. A scope reduction cannot meet that increase-in-scope test. For a modification that is not separate, paragraph 45 requires revised contract allocation, term and discount rate at the effective date. Paragraph 46(a) requires right-of-use derecognition and the gain/loss for a partial or full termination; the remaining liability remeasurement adjusts the surviving right-of-use asset. A mere decrease in price with unchanged scope belongs to paragraph 46(b), not automatic derecognition.

Represent the terminated rights explicitly, including the physical/use/time units and their carrying-cost allocation. In a homogeneous space lease with an unchanged remaining term and proportional rights, let fraction q of the rights be surrendered. The illustrative proportional allocation removes \(qA\) of the right-of-use carrying amount A and \(qL\) of the pre-modification liability L; the termination gain is \(q(L-A)\). This allocation requires the homogeneous-rights facts. Do not use the payment reduction as q or apply floor-area proportions to heterogeneous equipment or a pure term change without a supported allocation.

For \(A=80\) and \(L=100\), surrender 25\% of identical space with no termination fee. Remove right-of-use 20 and liability 25; recognize gain 5. The remaining \(A=60\) and \(L=75\) are then remeasured. Suppose the remaining contract requires one payment 70.4 after one year and the newly supported discount rate is 10\%. Its present value is 64. Reduce the remaining liability by 11 and the asset by 11, giving asset 49 and liability 64. The complete journal debits liability 36, credits the asset 31 and credits gain 5. It is wrong to send the entire liability reduction 36 to income or to omit the termination gain by adjusting the asset 36.

With opening cash 120/ROU 80/liability 100/equity 100, the event leaves cash 120, ROU 49/liability 64/equity 105. Over the one remaining year, zero residual and straight-line use give depreciation 49, interest 6.4 and payment 70.4. Closing cash and equity are 49.6, with no lease balances. Present-value calculation and the right-of-use consumption pattern are independent assumptions.

A separate fee of 2 per actual unit of usage, not index/rate dependent, not in-substance fixed and not already in the liability, is recognized under IFRS 16.38(b) when the condition occurs. Three units create expense 6/payable 6; payment reduces cash/payable 6. It is not a lease-liability remeasurement. Qualifying production use can instead include that cost in another asset under the applicable standard. Missing usage evidence, uncertain payment character, heterogeneous surrendered rights or a termination penalty require their own measurement before the conditional example can be applied.

\hypertarget{curtailment-settlement-and-the-asset-ceiling}{%
\subsubsection{Curtailment, settlement and the asset ceiling}\label{curtailment-settlement-and-the-asset-ceiling}}

The benefits record separates obligation D, fair-valued plan assets P, asset- ceiling adjustment, recognized net liability/asset, actuarial assumptions, remaining-year service cost, discount rate, legal discharge and actual payments. For a liability \(D-P\) is positive. If there is a surplus, the recognized asset is limited to the admissible economic benefit C; the added ceiling adjustment is \(\max(P-D-C,0)\). This simple representation assumes no further minimum-funding adjustment; any applicable IFRIC 14 obligation remains a separate admission dependency. The example supplies C, not a method for inferring legal recoverability.

IAS~19.99 \citep[IAS~19.99--112,120--126]{eu20231803} requires current assumptions and asset fair values immediately before and after amendment, curtailment or settlement. First recognize pre-event remeasurements, then the change in benefits/covered employees, then any later settlement. Past service cost arises from an amendment or curtailment and is recognized at the earlier date specified by paragraph 103; a related restructuring or termination-benefit recognition can therefore precede the formal change. The examples have no such earlier trigger. A settlement discharges obligations; ordinary scheduled benefits included in assumptions are not a settlement gain event. Buying insurance while retaining an obligation is also not settlement under paragraphs 111--112.

Measure past service and settlement gain/loss without the asset ceiling, then remeasure the ceiling under paragraph 101A. Service cost and settlement enter profit unless eligible capitalization applies; actuarial and ceiling remeasurements enter OCI under paragraph 120 and are not later recycled under paragraph 122. For the rest of the year, use updated service assumptions and the post-event net position/discount rate, allowing for subsequent contributions and benefits, under paragraphs 122A--126. An unchanged beginning-of-year interest base would miss this required reset.

Consider opening \(D=1{,}000\), \(P=900\), net liability 100 and company cash 500/equity 400. At the event date, current measurement gives \(D=1{,}100\), \(P=920\): actuarial loss 100 less asset gain 20 creates OCI loss 80 and liability 180. A curtailment reduces the freshly measured obligation by 150 to 950. Recognize past-service gain 150, leaving net liability 30. A later transaction fully discharges a measured obligation 200 for plan assets 170 plus company cash 20. Settlement price 190 gives gain 10; D falls to 750, P to 750, liability to zero and cash to 480. Profit 160 and OCI−80 reconcile equity from 400 to 480.

For the remaining half-year, assume current service cost 30, annual discount rate 5\%, simple half-year interest accrual and no contributions or benefit payments. Gross obligation and asset interest are both 18.75, so net interest is zero. Closing \(D=798.75\), \(P=768.75\), net liability 30, cash 480 and equity 450. Full period profit 130 and OCI−80 reconcile the change 50. Using the original net liability 100 for half-year interest would incorrectly add 2.5 expense. The simple accrual convention is an illustrative assumption, not a universal actuarial compounding rule.

To see the ceiling ordering, take D 100/P 130 and recoverable economic benefit 10. The surplus 30 yields recognized asset 10 and ceiling adjustment 20. An amendment reduces D by 20 to 80, while current evidence changes C to 15. Recognize past-service gain 20 without the ceiling. The new surplus 50 has ceiling adjustment 35, so its increase 15 is an OCI loss. The net recognized asset rises only 5 to 15. Recording just gain 5 would misclassify profit and OCI even though closing equity agrees.

\hypertarget{multi-employer-information-is-a-classification-decision}{%
\subsubsection{Multi-employer information is a classification decision}\label{multi-employer-information-is-a-classification-decision}}

IAS~19.32--36 \citep[IAS~19.32--41]{eu20231803} requires classification from actual plan terms, including constructive obligations. With adequate reliable allocation, a participant in a defined-benefit multi-employer plan recognizes its share of obligation, assets and cost. If the specific paragraph 34 information condition holds, account as if defined contribution and retain the paragraph 148 disclosures. Insufficient information is not proof that the plan economically has no defined-benefit exposure.

Paragraph 37 separately requires any contractual surplus/deficit-funding asset or liability and resulting profit effect. A synthetic agreed deficit payment 121 due in two years, with a supplied appropriate discount rate 10\%, has liability 100 and expense 100 initially, followed by interest 10 and 11 and settlement 121. Ordinary current-service contributions are distinct. Do not omit the deficit contract because ongoing contributions use defined-contribution accounting or recognize the nominal 121 immediately while also accruing interest.

Withdrawal or wind-up liabilities follow IAS 37 under paragraph 39. A plan among entities under common control is not a multi-employer plan under paragraph 40; paragraph 41's plan-wide information and separate-entity charge/allocation rules must be applied instead. The record must therefore retain employer relationships, allocation reliability, contributions, deficit agreements, withdrawal terms and required disclosures. Unknown relationship or contractual terms block the classification, not the preservation of the underlying evidence.
